% Ch6 WS2 -- Hess's law. Block 3.
\documentclass{apchem-worksheet}

\apcunitword{Chapter}
\apcunit{6}
\apcunittitle{Thermochemistry}
\apcdoctitle{Worksheet 2 \textbullet\ Hess's Law}
\apcsubtitle{Zumdahl \S6.3 \textbullet\ write the target first, then work backwards}

\begin{document}
\wsheader[For every problem, show which given equation you reversed or
scaled and by what factor. An answer with no manipulation shown cannot be
checked --- or given partial credit.]

\problem[]{State the two rules for manipulating a thermochemical equation,
and give the physical reason for each.
\answerlines[3]{Reversing a reaction reverses the sign of $\Delta H$,
because the sign records the direction of heat flow. Multiplying the
coefficients multiplies $\Delta H$ by the same factor, because enthalpy is
extensive --- twice the reaction transfers twice the heat.}}

\problem[]{Given \ce{N2(g) + 3 H2(g) -> 2 NH3(g)},
$\Delta H = \SI{-92.2}{\kilo\joule}$, find $\Delta H$ for:}
\begin{parts}
  \item \ce{2 NH3(g) -> N2(g) + 3 H2(g)}:
        \answerblank[3.2cm]{\SI{+92.2}{\kilo\joule}}
  \item \ce{1/2 N2(g) + 3/2 H2(g) -> NH3(g)}:
        \answerblank[3.2cm]{\SI{-46.1}{\kilo\joule}}
  \item \ce{2 N2(g) + 6 H2(g) -> 4 NH3(g)}:
        \answerblank[3.2cm]{\SI{-184.4}{\kilo\joule}}
\end{parts}

\problem[]{Verify Hess's law using Zumdahl's nitrogen oxides. Given
\begin{align*}
  \ce{N2(g) + O2(g) &-> 2 NO(g)}   &\Delta H &= \SI{180}{\kilo\joule}\\
  \ce{2 NO(g) + O2(g) &-> 2 NO2(g)} &\Delta H &= \SI{-112}{\kilo\joule}
\end{align*}}
\begin{parts}
  \item Add the two equations and write the net reaction.
        \answerlines[2]{\ce{N2(g) + 2 O2(g) -> 2 NO2(g)} --- the \ce{2 NO}
        cancels.}
  \item Compute $\Delta H$ for the net reaction.
        \answerblank[3cm]{\SI{+68}{\kilo\joule}}
  \item What property of enthalpy guarantees this equals the directly
        measured value? \answerblank[4cm]{it is a state function}
\end{parts}

\problem[]{Determine $\Delta H$ for \ce{C(s) + 1/2 O2(g) -> CO(g)} given
\begin{align*}
  \text{(i)}\ \ce{C(s) + O2(g) &-> CO2(g)}   &\Delta H &= \SI{-393.5}{\kilo\joule}\\
  \text{(ii)}\ \ce{2 CO(g) + O2(g) &-> 2 CO2(g)} &\Delta H &= \SI{-566.0}{\kilo\joule}
\end{align*}
\workspace[3.4cm]
\keyonly{\begin{apcanswer}Reverse and halve (ii):
\ce{CO2 -> CO + 1/2 O2}, $\Delta H = \SI{+283.0}{\kilo\joule}$. Add to (i);
\ce{CO2} cancels. $\Delta H = -393.5 + 283.0 =
\SI{-110.5}{\kilo\joule}$\end{apcanswer}}}

\problem[]{Determine $\Delta H$ for
\ce{2 C(s) + H2(g) -> C2H2(g)} given
\begin{align*}
  \text{(i)}\ \ce{C2H2(g) + 5/2 O2(g) &-> 2 CO2(g) + H2O(l)} &\Delta H &= \SI{-1300.}{\kilo\joule}\\
  \text{(ii)}\ \ce{C(s) + O2(g) &-> CO2(g)} &\Delta H &= \SI{-393.5}{\kilo\joule}\\
  \text{(iii)}\ \ce{H2(g) + 1/2 O2(g) &-> H2O(l)} &\Delta H &= \SI{-286}{\kilo\joule}
\end{align*}
\workspace[3.8cm]
\keyonly{\begin{apcanswer}Reverse (i): $+1300.$; take $2\times$(ii):
$-787.0$; take (iii) as written: $-286$. Sum:
$\Delta H = \SI{+227}{\kilo\joule}$ --- acetylene is endothermic to
form.\end{apcanswer}}}

\problem[]{Determine $\Delta H$ for \ce{N2H4(l) + O2(g) -> N2(g) + 2 H2O(l)}
given
\begin{align*}
  \text{(i)}\ \ce{N2H4(l) + O2(g) &-> N2(g) + 2 H2O(g)} &\Delta H &= \SI{-534}{\kilo\joule}\\
  \text{(ii)}\ \ce{H2O(l) &-> H2O(g)} &\Delta H &= \SI{+44}{\kilo\joule}
\end{align*}
\workspace[3cm]
\keyonly{\begin{apcanswer}Reverse (ii) and double it:
\ce{2 H2O(g) -> 2 H2O(l)}, $\Delta H = \SI{-88}{\kilo\joule}$. Add to (i):
$\Delta H = -534 - 88 = \SI{-622}{\kilo\joule}$\end{apcanswer}}}

\problem[]{A student solving problem 6 forgets to double the reversed step
and reports \SI{-578}{\kilo\joule}. Explain why the doubling is required.
\answerlines[2]{The target produces two moles of liquid water, so two moles
of vapour must be condensed --- and enthalpy is extensive, so the
condensation enthalpy must be doubled along with the coefficients.}}

\begin{spiralstrip}{Chapter 6 \textbullet\ blocks 1--2}
  \item $\Delta E$ if $q = +300$ J and $w = -100$ J:
        \answerblank[2.6cm]{$+200$ J}
  \item $\Delta H = q$ under what condition?
        \answerblank[3.4cm]{constant pressure}
  \item $q$ for \SI{100.}{\gram} water warmed \SI{5.00}{\celsius}:
        \answerblank[3cm]{\SI{2092}{\joule}}
  \item Specific heat of water:
        \answerblank[3.6cm]{\SI{4.184}{\joule\per\gram\per\celsius}}
  \item Sign of $w$ when a gas is compressed:
        \answerblank[2.2cm]{positive}
\end{spiralstrip}

\end{document}
